% =======================================================
% 8. PREVENTIVO A FINIRE (PaF)
% =======================================================
\section{Preventivo a Finire (PaF)}
\label{sec:paf}

Il Preventivo a Finire viene aggiornato al termine di ogni sprint a valle del ciclo di controllo del Piano di Progetto. L'aggiornamento tiene conto del consuntivo di periodo, della Sprint Retrospective\textsubscript{G}, delle misure correttive individuate, dell'aggiornamento dei rischi e della revisione della pianificazione futura.

Il suo scopo è fornire una stima aggiornata del costo complessivo del progetto, tenendo conto del costo già sostenuto, delle risorse residue e del costo stimato per completare le attività rimanenti.

\subsection{Analisi scostamenti}

Il monitoraggio delle ore residue per ruolo, aggiornato al termine dello
Sprint 19, evidenzia una distribuzione non uniforme degli scostamenti rispetto
alla preventivazione iniziale:

\begin{table}[H]
\centering
\renewcommand{\arraystretch}{1.3}
\begin{tabular}{l r r}
\toprule
\textbf{Ruolo} & \textbf{Ore residue} & \textbf{Costo residuo} \\
\midrule
Responsabile   & 23{,}33 & 700{,}00~\euro \\
Amministratore & \color{red}{-43{,}83} & \color{red}{-876{,}67~\euro} \\
Analista       & 54{,}67 & 1.366{,}67~\euro \\
Progettista    & 98{,}17 & 2.454{,}17~\euro \\
Programmatore  & 42{,}00 & 630{,}00~\euro \\
Verificatore   & 51{,}17 & 767{,}50~\euro \\
\midrule
\textbf{Totale} & \textbf{225{,}51} & \textbf{5.041{,}67~\euro} \\
\bottomrule
\end{tabular}
\caption{Ore e costo residuo per ruolo dopo lo Sprint 19}
\end{table}

Il ruolo di Amministratore\textsubscript{G} è l'unico a presentare un residuo
negativo (-43,83 ore, pari a -876,67~\euro): il gruppo ha impiegato più ore di
quelle preventivate per mantenere aggiornati gli strumenti di gestione, la
documentazione di progetto, il tracciamento delle attività e i consuntivi di
sprint. Questo scostamento, già rilevato in fase RTB\textsubscript{G} (Sez.
RR/RO6 dell'Analisi dei Rischi), è di origine prevalentemente RTB: nel periodo
di implementazione dell'MVP (Sprint 12 -- Sprint 19), la somma dei
$\Delta$~ore di Amministratore sui singoli sprint è pari a -0,01 ore, cioè
sostanzialmente nulla. Il deficit non si è quindi aggravato durante la PB, ma
nemmeno riassorbito: resta un debito ereditato dalla RTB, non un problema
ricorrente della fase corrente.

Nello stesso periodo (Sprint 12 -- Sprint 19) emergono due scostamenti minori
e coerenti con la natura della fase: Progettista e Programmatore registrano
entrambi un $\Delta$ cumulato di +1,00 ora sugli otto sprint, riconducibile
all'intensificarsi delle attività di progettazione architetturale e codifica
dell'MVP; Verificatore registra un $\Delta$ cumulato di +0,16 ore, segno che
il carico aggiuntivo dovuto all'introduzione della suite di test automatizzati
e della pipeline di CI\textsubscript{G} è stato assorbito senza scostamenti
significativi. Analista e Responsabile risultano invece sostanzialmente in
linea con il preventivo nello stesso periodo ($\Delta$ cumulato pari a 0,00
ore per entrambi), poiché le attività di analisi si sono consolidate dopo la
validazione RTB dell'Analisi dei Requisiti.

Il residuo complessivo (225,51 ore, pari a 5.041,67~\euro) resta ampiamente
positivo: il deficit di Amministratore viene compensato dai residui positivi
degli altri ruoli, in particolare Progettista e Analista, senza che ciò
comporti un superamento del budget economico complessivo preventivato in fase
di candidatura.

\subsection{Ripianificazione}

A seguito degli scostamenti descritti, il gruppo ha ribilanciato nel corso
della PB\textsubscript{G} la distribuzione delle ore residue tra ruoli,
mantenendo invariato il vincolo economico complessivo definito nel preventivo
iniziale. In particolare:
\begin{itemize}
    \item il debito di ore di Amministratore ereditato dalla RTB\textsubscript{G}
    non è stato ricostituito durante la PB (delta netto sostanzialmente nullo
    su Sprint 12 -- Sprint 19), grazie a una più attenta pianificazione delle
    attività amministrative in sede di Sprint Planning\textsubscript{G};
    \item le ore residue del ruolo di Analista, resesi disponibili dal
    consolidamento dell'Analisi dei Requisiti dopo la validazione RTB, sono
    state impiegate nella pianificazione delle attività rimanenti, con
    particolare attenzione alla verifica e all'aggiornamento della
    documentazione tecnica per la PB (Specifica Tecnica, Manuale Utente);
    \item i lievi sovraccarichi di Progettista e Programmatore, coerenti con
    l'intensità delle attività di progettazione e codifica dell'MVP\textsubscript{G},
    sono stati assorbiti dal residuo positivo accumulato su questi ruoli
    (rispettivamente 98,17 e 42,00 ore residue a fine Sprint 19) senza
    necessità di ulteriori interventi correttivi.
\end{itemize}

Per gli sprint rimanenti fino alla Revisione PB (2026-09-10), il gruppo
continua a monitorare le ore residue per ruolo al termine di ogni sprint,
confermando il budget iniziale come vincolo massimo non superabile. Il
residuo disponibile (225,51 ore, 5.041,67~\euro) è ritenuto sufficiente per
completare le attività pianificate: redazione finale della documentazione
di PB, eventuali correzioni emerse dalla Validazione MVP\textsubscript{G} con
la proponente e preparazione della Revisione PB, senza necessità di
un'ulteriore ripianificazione strutturale.

Eventuali ulteriori scostamenti verranno comunque valutati nelle successive
retrospettive e recepiti nel Preventivo a Finire, aggiornando la stima delle
risorse residue senza modificare il perimetro economico complessivo del progetto.

Le azioni correttive adottate durante la PB, confermate efficaci dai dati di
Sprint 12 -- Sprint 19, sono:
\begin{itemize}
    \item maggiore attenzione alla pianificazione delle attività amministrative e
    di gestione durante lo Sprint Planning, evitando che vengano considerate
    attività accessorie o non preventivate;
    \item monitoraggio più frequente delle ore effettive tramite gli strumenti di
    consuntivazione, con verifica del residuo per ruolo al termine di ogni sprint;
    \item mantenimento del budget complessivo iniziale come vincolo massimo non
    superabile, anche in presenza di scostamenti localizzati su singoli ruoli.
\end{itemize}

Tali azioni risultano coerenti con le strategie di mitigazione definite nel
Piano di Progetto e hanno permesso di stabilizzare lo scostamento di
Amministratore durante la PB, pur senza riassorbire il debito ereditato dalla
RTB, preservando in ogni caso il rispetto delle scadenze e del budget
complessivo.

\subsection{Quadro economico aggiornato}

\begin{table}[H]
\centering
\renewcommand{\arraystretch}{1.3}
\begin{tabular}{l r}
\toprule
\textbf{Voce} & \textbf{Importo} \\
\midrule
Budget iniziale preventivato  & 11.665{,}00~\euro \\
Costo sostenuto finora        & 6.623{,}32~\euro \\
Costo stimato per completare  & 5.041{,}68~\euro \\
\midrule
\textbf{Nuovo totale stimato} & \textbf{11.665{,}00~\euro} \\
\bottomrule
\end{tabular}
\caption{Quadro economico aggiornato del Preventivo a Finire dopo lo Sprint 19}
\end{table}

Al termine dello Sprint 19 risultano sostenuti 6.623{,}32~\euro, pari a circa il 56{,}8\% del budget complessivo preventivato. Le risorse residue, pari a 5.041{,}68~\euro, sono ritenute sufficienti per completare le attività pianificate senza superare il costo massimo accettato in fase di candidatura.